

\section{Исследование и построение решения задачи}
\label{sec:Chapter3} \index{Chapter3}

В данной главе предлагается решение задачи, поставленной в разделе~\ref{sec:Chapter1}. Исходная задача~--- повышение качества многошаговых рассуждений LLM посредством использования явной графовой структуры~--- разбивается на ряд более мелких подзадач, для каждой из которых формулируется и обосновывается конкретное техническое решение. Декомпозиция выполняется в логическом порядке: от подготовки данных и формализации метрик к построению фреймворка, выбору базовой модели и проектированию функции графовой награды для дообучения с подкреплением.

\subsection{Декомпозиция задачи}
\label{subsec:decomposition}

Сформулированная в разделе~\ref{subsec:task-description} задача требует разработки полного пайплайна, включающего как генерацию данных, так и обучение модели. Для того чтобы свести её к набору технически реализуемых шагов, выделим следующие подзадачи:

\begin{enumerate}
    \setlength{\itemsep}{\smallskipamount}
    \item \textit{Построение датасета вопросов с эталонными подграфами рассуждения.} Необходимо сформировать набор пар $(q, G^*)$, где $q$~--- вопрос, требующий многошагового рассуждения, $G^*$~--- эталонный подграф, описывающий корректную цепочку вывода (раздел~\ref{subsec:dataset-construction}).
    \item \textit{Формализация метрик структурного соответствия.} Требуется определить количественные меры близости между сгенерированным графом $\hat{G}$ и эталонным графом $G^*$, пригодные как для оценки моделей, так и для использования в качестве сигнала обучения (раздел~\ref{subsec:metrics-formalization}).
    \item \textit{Извлечение графа рассуждения из текстовой генерации.} Поскольку LLM выдаёт ответ в текстовой форме, необходима процедура автоматического преобразования текста в множество пар <<посылка~--- следствие>>, образующих граф $\hat{G}$ (раздел~\ref{subsec:graph-extraction}).
    \item \textit{Разработка фреймворка smart-thinking-llm.} Все перечисленные компоненты должны быть объединены в модульную программную систему, поддерживающую воспроизводимые эксперименты с различными моделями, датасетами и алгоритмами обучения (раздел~\ref{subsec:framework-architecture}).
    \item \textit{Выбор базовой модели для дообучения.} На основе сравнительной оценки нескольких открытых LLM по графовым метрикам необходимо обосновать выбор архитектуры, демонстрирующей наилучший баланс качества рассуждений и доступности для дообучения с подкреплением (раздел~\ref{subsec:model-selection}).
    \item \textit{Проектирование функции графовой награды.} Требуется построить функцию награды для RL-этапа, объединяющую корректность финального ответа, соблюдение формата и структурное соответствие сгенерированного графа эталонному (раздел~\ref{subsec:reward-design}).
    \item \textit{Алгоритм обучения с подкреплением.} Подзадача состоит в выборе и настройке алгоритма RL и соответствующих программных платформ, обеспечивающих устойчивое обучение модели на сформированных данных и с разработанной функцией награды (раздел~\ref{subsec:training-algorithm}).
\end{enumerate}

В разделе~\ref{subsec:overall-pipeline} приведена общая схема пайплайна, объединяющая решения всех подзадач.

\subsection{Построение датасета с эталонными подграфами рассуждения}
\label{subsec:dataset-construction}

Для обучения и оценки модели требуются пары вида $(q, G^*)$. В работе используются два независимых источника таких данных, что обеспечивает диверсификацию обучающего сигнала и проверку обобщающей способности подхода.

\textbf{RuleTaker.} Датасет RuleTaker~\cite{clark2020ruletaker} содержит синтетические задачи логического вывода: задаётся набор атомарных фактов и правил <<если~\ldots, то~\ldots>>, а вопрос требует определения, следует ли заданное утверждение из имеющейся базы знаний (метка \texttt{entailment} или \texttt{not entailment}). Для каждой задачи известна минимальная цепочка вывода, которая представляется в виде ориентированного графа: вершинами выступают высказывания (факты, условия правил и промежуточные следствия), а рёбрами~--- отношения логического следования <<посылка $\rightarrow$ заключение>>.

\textit{Предобработка и фильтрация.} Поставляемый RuleTaker граф рассуждения содержит вершины трёх типов~--- \texttt{QUESTION}, \texttt{CONDITION} и \texttt{RULE}, причём вершины \texttt{RULE} выступают <<мостами>> между условиями и следствиями и не несут содержательного шага рассуждения. Поэтому реализована процедура схлопывания (collapse) вершин типа \texttt{RULE}: для каждой тройки \texttt{CONDITION/QUESTION} $\rightarrow$ \texttt{RULE} $\rightarrow$ \texttt{CONDITION/QUESTION} вершина \texttt{RULE} удаляется, а её входящие и исходящие рёбра объединяются в прямое ребро. Для каждого заключения возможны несколько альтернативных наборов посылок, что представляется отображением вида $\{\text{conclusion} \mapsto [[\text{cond}_1, \text{cond}_2], [\text{cond}_3]]\}$, где каждый внутренний список~--- одно из допустимых обоснований данного заключения.

Дополнительно из исходного RuleTaker отбираются только конфигурации, содержащие отрицания в отношениях (наличие токена \texttt{not} в типе ребра) и относящиеся к подмножеству \texttt{NatLang}~--- задачам, сформулированным на естественном языке. Полученный датасет опубликован под именем \texttt{KGReasoning} и разбивается на обучающую и валидационную части в соотношении $0{.}9{:}0{.}1$ с фиксированным значением \texttt{seed}~$=42$.

Каждый пример датасета хранится в виде структуры с полями \texttt{prompt} (системный и пользовательский промпты в формате chat), \texttt{label} (эталонная метка \texttt{entailment}/\texttt{not entailment}) и \texttt{proof\_rules} (отображение $\Pi^*$ из определения~\ref{def:proof-graph}), а также служебных полей \texttt{context}/\texttt{question} в \texttt{extra\_info}. Отбор примеров с отрицаниями повышает логическую сложность: задача перестаёт сводиться к поиску прямой цепочки фактов и требует учёта опровергающих условий. Валидационная оценка проводится на срезе до $10\,000$ примеров (раздел~\ref{subsec:impl-results}).

\textbf{Wikidata.} Второй источник данных формируется на основе графа знаний Wikidata~\cite{vrandecic2014wikidata} и его текстовой проекции Wikidata5m. Процедура построения датасета включает следующие шаги:
\begin{enumerate}
    \setlength{\itemsep}{\smallskipamount}
    \item Из графа Wikidata случайным образом выбирается стартовая сущность $v_0 \in V$.
    \item От вершины $v_0$ строится путь длины $k \in \{1, 2, 3\}$ с фильтрацией редких и низкокачественных типов отношений; реализованы отдельные поисковики путей $1$-, $2$- и $3$-hop, а также композиционных $3{+}2$ путей.
    \item Полученная последовательность троек $(v_0, r_1, v_1), \ldots, (v_{k-1}, r_k, v_k)$ образует эталонный подграф $G^*$.
    \item На основе подграфа с помощью внешней мощной LLM формулируется вопрос $q$ на естественном языке, ответом на который является финальная сущность $v_k$. В промпте требуется, чтобы ответ не мог быть получен без последовательного применения всех рёбер пути.
    \item Полученная пара $(q, G^*)$ валидируется: вопрос подаётся на вход другой LLM в режиме без доступа к графу, и принимается в датасет только в случае, если ответ не совпадает с эталонным.
\end{enumerate}

\textit{Стратегии поиска путей и генерации вопросов.} Для повышения разнообразия обучающего сигнала реализовано несколько типов поисковиков структур и соответствующих генераторов вопросов:
\begin{itemize}
    \setlength{\itemsep}{\smallskipamount}
    \item \textit{Прямые пути} длины $1$, $2$ и $3$ (\texttt{two\_hop\_finder}, \texttt{three\_hop\_finder})~--- линейная цепочка отношений, в которой ответ требует последовательного применения всех рёбер;
    \item \textit{Bridge-конфигурации} ($2{\to}1$ и $3{\to}1$, генераторы \texttt{bridge\_21\_generator}, \texttt{bridge\_31\_generator})~--- несколько путей, сходящихся в общей сущности-мосте, что требует объединения фактов;
    \item \textit{Композиционные пути} $3{+}2$ (\texttt{compositional\_3\_2\_finder}, \texttt{compositional\_question\_generator})~--- комбинация двух подпутей разной длины, повышающая глубину рассуждения.
\end{itemize}
Поиск сопровождается фильтрацией редких и низкокачественных типов отношений, а вопросы формулируются внешней LLM по соответствующему шаблону промпта (подкаталог \texttt{prompts/}).

Пример двухшагового Wikidata-пути приведён на рис.~\ref{fig:wikidata-path}: вопрос <<\textit{What is the top-level Internet domain for the country where Miyankuh-e Gharbi is located?}>> требует последовательного применения двух отношений (\texttt{located in} и \texttt{top-level Internet domain}), и ответ \texttt{.ir} не может быть получен без прохождения обоих рёбер.

\begin{figure}[H]
\centering
\begin{tikzpicture}[
    ent/.style={draw, rounded corners, align=center, font=\footnotesize, minimum height=8mm, inner sep=4pt, fill=black!4},
    arr/.style={-{Stealth[length=2mm]}, thick},
]
\node[ent] (v0) {Miyankuh-e Gharbi\\\scriptsize $v_0$};
\node[ent, right=26mm of v0] (v1) {Iran\\\scriptsize $v_1$};
\node[ent, right=26mm of v1] (v2) {.ir\\\scriptsize $v_2$ (ответ)};
\draw[arr] (v0) -- node[above, font=\scriptsize]{located in} (v1);
\draw[arr] (v1) -- node[above, font=\scriptsize]{ccTLD} (v2);
\end{tikzpicture}
\caption{Пример эталонного Wikidata-пути длины $k=2$}
\label{fig:wikidata-path}
\end{figure}

Использование двух качественно различных источников~--- синтетических логических задач (RuleTaker) и фактических многошаговых вопросов (Wikidata)~--- позволяет независимо оценить способность модели к формально-логическому и к фактическому многошаговому рассуждению.

\subsection{Формализация метрик структурного соответствия}
\label{subsec:metrics-formalization}

Для пары графов $(\hat{G}, G^*)$, где $\hat{G} = (\hat{V}, \hat{E})$~--- сгенерированный граф, $G^* = (V^*, E^*)$~--- эталонный, определяются метрики на уровне вершин, рёбер и структуры графа в целом.

\textbf{Метрики на уровне вершин:}
\begin{align}
    P_V = \frac{|\hat{V} \cap V^*|}{|\hat{V}|}, \qquad
    R_V = \frac{|\hat{V} \cap V^*|}{|V^*|}, \qquad
    F_{1,V} = \frac{2 P_V R_V}{P_V + R_V}.
\end{align}

\textbf{Метрики на уровне рёбер.} Аналогично определяются Edge Precision $P_E$, Edge Recall $R_E$ и Edge F1 $F_{1,E}$ с заменой множеств $\hat{V}, V^*$ на $\hat{E}, E^*$. Сравнение рёбер выполняется по тройкам $(s, r, o)$; для RuleTaker, где ребро представляет отношение <<посылка $\rightarrow$ заключение>>, сравниваются упорядоченные пары вершин.

\textbf{Метрики на уровне структуры.} Дополнительно используется Graph Edit Distance (GED)~--- минимальное число операций вставки, удаления и замены вершин и рёбер, необходимых для преобразования $\hat{G}$ в $G^*$. На её основе вводится нормированная мера Graph Edit Similarity:
\begin{align}
    \mathrm{GES}(\hat{G}, G^*) = 1 - \frac{\mathrm{GED}(\hat{G}, G^*)}{\max(|\hat{V}| + |\hat{E}|,\; |V^*| + |E^*|)} \in [0, 1].
    \label{eq:ges}
\end{align}
Значение $\mathrm{GES} = 1$ соответствует изоморфным графам, $\mathrm{GES} = 0$~--- полностью различным. Вычисление GED для произвольных графов NP-трудно, однако в условиях данной задачи (графы из единиц~--- десятков вершин) применима реализация на основе библиотеки \texttt{networkx}.

\textbf{Покрытие как сигнал обучения: рекурсивная валидация подграфа.} Перечисленные выше метрики оценивают граф <<как множество>> и не учитывают структурной согласованности шагов рассуждения. Для использования в качестве сигнала RL-обучения разработана отдельная процедура валидации, явно проверяющая корректность каждого шага вывода.

Пусть $\hat{G}$~--- сгенерированный граф, представленный множеством рёбер $(p \rightarrow c)$, где $p$~--- посылка, $c$~--- заключение, а $G^*$ задан как отображение
\[
\Pi^* : c \mapsto \{S_1, S_2, \ldots\}, \qquad S_i = \{p^{(i)}_1, p^{(i)}_2, \ldots\},
\]
сопоставляющее каждому заключению список допустимых наборов посылок. Процедура \texttt{validate\_subgraph} устроена следующим образом: обход стартует от вершины-вопроса $q$ и для каждого узла-заключения $c$, входящего в $\hat{G}$, собирается множество посылок $\hat{S}(c) \subseteq \hat{V}$, использованных моделью. Узел $c$ считается \emph{полностью валидным} ($\mathrm{score}(c) = 1$), если существует $S_i \in \Pi^*(c)$ такой, что $\hat{S}(c) = S_i$, \emph{частично валидным} ($\mathrm{score}(c) = 0{.}5$), если такого совпадения нет, но $c$ присутствует в $G^*$, и \emph{невалидным} ($\mathrm{score}(c) = 0$) в противном случае. Покрытие графа определяется как
\begin{align}
    \mathrm{Cov}(\hat{G}, G^*) = \frac{1}{|\hat{V}_c|} \sum_{c \in \hat{V}_c} \mathrm{score}(c) \in [0, 1],
    \label{eq:coverage}
\end{align}
где $\hat{V}_c$~--- множество всех узлов-заключений в $\hat{G}$. Отдельно предусмотрен случай отрицательного примера: если $a^* = \texttt{not entailment}$ и модель сгенерировала пустой граф (что соответствует корректному отказу от вывода), полагается $\mathrm{Cov} = 1$.

В отличие от агрегатов $P_V/R_V/P_E/R_E$, формула~(\ref{eq:coverage}) учитывает, что каждое заключение должно быть выведено из \emph{полного и корректного} множества посылок, а не просто содержать пересечение с эталонными вершинами и рёбрами по отдельности. Это делает $\mathrm{Cov}$ существенно более информативным сигналом для обучения политики.

\textbf{Сопоставление метрик как обучающего сигнала.} Три семейства метрик различаются по пригодности для использования в качестве награды. Метрики на множествах ($P/R/F_1$) дёшевы, но оценивают граф как неупорядоченное множество элементов и не штрафуют структурную несвязность: ответ, содержащий правильные вершины и рёбра в некорректной комбинации, получает завышенную оценку. Мера GES учитывает структуру целиком, но её вычисление дорого (NP-трудная задача), а сама величина слабо градуирована для частично корректных цепочек и плохо локализует ошибку. Покрытие $\mathrm{Cov}$ сочетает достоинства обоих: вычисляется за линейное время, штрафует неполные и некорректные наборы посылок и даёт промежуточные градации (полная/частичная валидность узла). Сводное сравнение приведено в таблице~\ref{tab:metric-signal}.

\begin{table}[H]
\centering
\caption{Пригодность метрик в качестве сигнала RL-обучения}
\label{tab:metric-signal}
\small
\begin{tabular}{|L{3.6cm}|C{3.0cm}|C{2.4cm}|C{3.4cm}|}
\hline
\textbf{Свойство} & \textbf{$P/R/F_1$} & \textbf{GES} & \textbf{Покрытие $\mathrm{Cov}$} \\
\hline
Сложность вычисления & линейная & экспон. (NP) & линейная \\
\hline
Учёт структуры вывода & нет & да & да (по узлам) \\
\hline
Градации частичной корректности & слабые & слабые & да ($1/0{,}5/0$) \\
\hline
Локализация ошибки & нет & нет & да (узел) \\
\hline
Пригодность как награда & низкая & средняя & высокая \\
\hline
\end{tabular}
\end{table}

Поэтому в качестве графовой компоненты награды (раздел~\ref{subsec:reward-design}) выбрано именно покрытие $\mathrm{Cov}$, тогда как $F_1$ и GES применяются преимущественно для итоговой оценки моделей.

\subsection{Извлечение графа рассуждения из текстовой генерации}
\label{subsec:graph-extraction}

Поскольку LLM производит ответ в свободной текстовой форме, необходима процедура автоматического преобразования текста в граф $\hat{G}$. Были разработаны два подхода, применяемые в разных режимах работы фреймворка.

\textbf{Подход 1: специализированный формат ответа и JSON-разбор.} Этот подход является основным для обучаемой модели и используется как при сэмплировании RL-кандидатов, так и при инференсе обученной политики. Системный промпт задаёт жёсткий формат вывода, явно разделяющий внутренние размышления и финальный структурированный ответ:
\begin{center}
\texttt{<think>...свободные размышления...</think>}\\
\texttt{<answer>[[evidence, conclusion], ...]$\backslash$nentailment | not entailment</answer>}
\end{center}
Внутри блока \texttt{<answer>} модель обязана выдать (i)~JSON-массив рёбер графа рассуждения, каждое из которых представлено парой <<строка-посылка, строка-заключение>>, и (ii)~на отдельной строке~--- финальный ответ задачи (для RuleTaker~--- бинарную метку). Извлечение графа реализуется в виде последовательности: парсинг тегов \texttt{<think>/<answer>}, извлечение JSON-массива регулярным выражением, его разбор стандартным парсером и построение графа. Метка ответа нормализуется через словарь алиасов (\texttt{true/yes/follows} $\to$ \texttt{entailment}; \texttt{false/no/contradiction/does not follow} $\to$ \texttt{not entailment}). Достоинство подхода~--- скорость, детерминированность и возможность точечной диагностики (каждый этап парсинга даёт собственный сигнал, см. раздел~\ref{subsec:reward-design}); недостаток~--- требует, чтобы модель устойчиво поддерживала формат, что обеспечивается за счёт включения штрафов за нарушение формата непосредственно в функцию награды.

\textbf{Подход 2: LLM-ассесор со структурированным выходом.} Для оценки моделей, не приученных к строгому формату (например, базовых open-source LLM на этапе сравнения в разделе~\ref{subsec:model-selection}), а также для построения референсных графов на этапе подготовки данных используется отдельный модуль извлечения~--- \texttt{StructuredGraphCreator}. Текстовый ответ модели подаётся на вход вспомогательной LLM (использовались \texttt{Qwen/Qwen3-8B} и \texttt{meta-llama/Llama-3.1-8B-Instruct}); запрос оформляется в виде Structured Output по схеме Pydantic (классы \texttt{RuleTakerSimplePath}, \texttt{ReasoningGraph}, \texttt{ReasoningGraphNodes}, \texttt{ReasoningGraphConnections}, \texttt{ReasoningGraphIdBased}), что гарантирует синтаксически корректный JSON-выход. Для задач RuleTaker реализованы также варианты двухпроходного извлечения (\texttt{GraphCreatorWithLLMSelector}) и подхода с эмбеддингами для сопоставления естественноязыковых высказываний (\texttt{GraphCreatorWithEmbeddings}), позволяющего сводить семантически эквивалентные формулировки к одной вершине. Параметры запросов унифицированы (\texttt{temperature} = $0{.}2$, \texttt{max\_tokens} = $5000$) для воспроизводимости.

\textbf{Сравнение подходов и выбор.} На пилотной выборке оба подхода сопоставлены по доле успешно извлечённых эталонных троек. Подход~1 обеспечивает существенно большую пропускную способность и не требует вспомогательной модели, поэтому используется в основном RL-пайплайне; подход~2 применяется как fallback и для оценки моделей, не способных самостоятельно соблюдать формат, а также как генератор обучающих сигналов на этапе создания датасета. Сопоставление подходов по ключевым характеристикам приведено в таблице~\ref{tab:extraction}.

\begin{table}[H]
\centering
\caption{Сравнение подходов к извлечению графа рассуждения из текста}
\label{tab:extraction}
\small
\begin{tabular}{|L{4.2cm}|C{4.0cm}|C{4.0cm}|}
\hline
\textbf{Характеристика} & \textbf{Подход 1: JSON-разбор} & \textbf{Подход 2: LLM-ассесор} \\
\hline
Вспомогательная модель & не требуется & требуется (Qwen3-8B и др.) \\
\hline
Скорость & высокая & низкая (вызов LLM) \\
\hline
Детерминированность & да & нет \\
\hline
Требование к формату модели & строгое & отсутствует \\
\hline
Основное применение & RL-цикл, инференс & fallback, оценка, генерация данных \\
\hline
\end{tabular}
\end{table}

\subsection{Архитектура фреймворка smart-thinking-llm}
\label{subsec:framework-architecture}

Все компоненты предлагаемого решения объединены в программный фреймворк \texttt{smart-thinking-llm}, спроектированный по модульному принципу. Архитектура включает следующие подсистемы:

\begin{itemize}
    \setlength{\itemsep}{\smallskipamount}
    \item \textbf{Модуль данных} (\texttt{datasets/}): абстрактный базовый класс \texttt{BaseDataset}, конкретные реализации \texttt{RuleTakerDataset} и \texttt{WikidataDataset}, общие структуры данных (\texttt{data\_classes.py}).
    \item \textbf{Модуль поиска паттернов} (\texttt{find\_patterns/}): обход графа Wikidata для отбора путей с заданной структурой (\texttt{two\_hop\_finder.py}, \texttt{three\_hop\_finder.py}, \texttt{compositional\_3\_2\_finder.py}) и общий запускатель \texttt{pattern\_runner.py}.
    \item \textbf{Модуль генерации вопросов} (\texttt{generation/}): набор генераторов синтетических вопросов по найденным подграфам (\texttt{basic\_question\_generator.py}, \texttt{bridge\_21\_generator.py}, \texttt{bridge\_31\_generator.py}, \texttt{compositional\_question\_generator.py}) и \texttt{structured\_ruletaker\_runner.py} для генерации цепочек рассуждений в формате Structured Output.
    \item \textbf{Модуль шаблонов} (\texttt{prompts/}): системные промпты для генерации троек, отбора описаний и для разных режимов RuleTaker-промптинга (подкаталог \texttt{ruletaker\_structured/}).
    \item \textbf{Модуль извлечения графов} (\texttt{tools/graph\_creator/}): реализация подходов из раздела~\ref{subsec:graph-extraction}~--- \texttt{GraphCreator}, \texttt{StructuredGraphCreator}, \texttt{GraphCreatorWithLLMSelector}, \texttt{GraphCreatorWithEmbeddings}.
    \item \textbf{Модуль структурированного вывода} (\texttt{tools/structured\_outputs.py}): обёртки над OpenAI-совместимым API инференс-сервера для запросов с Pydantic-схемами.
    \item \textbf{Модуль метрик} (\texttt{tools/metrics/}): класс \texttt{RuleTakerMetricCalculator}, реализующий \texttt{calculate\_path\_metrics\_graph} (на основе GED), \texttt{calculate\_simple\_path\_metrics} (на множествах троек) и \texttt{calculate\_reasoning\_graph\_metrics}; вычисления опираются на \texttt{networkx}.
    \item \textbf{Модуль вспомогательных утилит} (\texttt{tools/}): конвертеры путей и формата RuleTaker (\texttt{path\_converter.py}, \texttt{ruletaker\_converter.py}), работа с графами (\texttt{graph.py}), разбиение файлов (\texttt{split\_file.py}).
    \item \textbf{Модуль логирования} (\texttt{tensorboard\_logger.py}): унифицированное логирование метрик и компонент награды в TensorBoard.
\end{itemize}

Само дообучение реализовано в репозитории-спутнике \texttt{TunePPO}, содержащем конфигурации алгоритма GRPO, реализацию функции графовой награды для RuleTaker (\texttt{verl\_exp/edge\_reasoning\_reward.py}) и многошаговой графовой награды для Wikidata (\texttt{GraphMultihopQAReward} в \texttt{ppotune/reward.py}), а также скрипты подготовки данных (\texttt{verl\_exp/ruletaker\_prepare\_data\_kg.py}, \texttt{verl\_exp/ruletaker\_prepare\_data\_struct.py}). Такое разделение обусловлено тем, что smart-thinking-llm отвечает за всё, что относится к данным, графам и инференсу, тогда как TunePPO~--- за RL-цикл поверх двух промышленных платформ обучения: VERL~\cite{sheng2024verl} и собственной реализации на базе torchtune. Модульность архитектуры обеспечивает выполнение нефункционального требования о возможности замены отдельных компонентов (раздел~\ref{subsec:requirements}): для проведения нового эксперимента достаточно реализовать новый класс модели, датасета, экстрактора графа или функции награды, не изменяя остальные части системы.

\subsection{Выбор базовой модели для дообучения}
\label{subsec:model-selection}

Для последующего дообучения с подкреплением необходимо выбрать базовую модель, удовлетворяющую двум условиям: (1)~демонстрировать измеримое, но неполное качество рассуждений~--- чтобы существовал потенциал для улучшения, и (2)~быть достаточно компактной для обучения на доступных вычислительных ресурсах (одна-две GPU класса H100/A100).

\textbf{Методология сравнительной оценки.} Каждая из рассматриваемых моделей запускалась на едином подмножестве датасета RuleTaker в режиме pass@1. Для моделей, не приученных к строгому формату ответа, использовалось извлечение графа подходом~2 из раздела~\ref{subsec:graph-extraction} (\texttt{StructuredGraphCreator}); для моделей семейства Qwen3, обученных следовать заданному формату \texttt{<think>...</think><answer>...</answer>}, применялся прямой парсинг JSON-блока из \texttt{<answer>}. Оценка проводилась по трём метрикам: точности предсказания метки (\texttt{entailment}/\texttt{not entailment}), $F_1$ по рёбрам реконструированного графа и интегральной метрике GES~(\ref{eq:ges}).

\textbf{Рассматриваемые модели.} Анализу подверглись следующие семейства открытых моделей: Qwen3 в конфигурациях 0.6B, 1.7B, 4B, 8B, 14B и 32B параметров; Llama-3.1-8B-Instruct; OLMo-2-Think 1B и 7B; GPT-OSS 20B и 120B. Выбор именно этих семейств обусловлен наличием в них режима явных рассуждений (\texttt{thinking mode}) и открытыми лицензиями, допускающими дообучение. Качественная сводка по группам моделей приведена в таблице~\ref{tab:model-tiers}.

\begin{table}[H]
\centering
\caption{Качественная оценка групп моделей-кандидатов для RL-дообучения}
\label{tab:model-tiers}
\small
\begin{tabular}{|L{3.4cm}|C{2.6cm}|C{2.8cm}|C{3.2cm}|}
\hline
\textbf{Группа} & \textbf{Соблюдение формата} & \textbf{Потенциал улучшения} & \textbf{Доступность для RL (1 GPU)} \\
\hline
$<$1B (Qwen3-0.6B, OLMo-1B) & неустойчивое & высокий & высокая \\
\hline
1.7B--4B (Qwen3) & уверенное & средний & высокая \\
\hline
7B--8B (Qwen3-8B, Llama-3.1-8B, OLMo-7B) & уверенное & низкий (насыщение) & ограниченная \\
\hline
$\geq$14B (Qwen3-14B/32B, GPT-OSS) & уверенное & очень низкий & требует мультиGPU \\
\hline
\end{tabular}
\end{table}

\textbf{Выводы и финальный выбор.} Модели объёмом 8B и выше демонстрируют качество, близкое к насыщению по точности предсказания метки, оставляя мало пространства для улучшения, и при этом требуют существенно больших ресурсов для LoRA-дообучения с RL. По мере уменьшения размера модели качество реконструкции графа снижается, однако тем самым возрастает потенциал его улучшения за счёт дообучения. Наибольший интерес с точки зрения поставленной задачи представляют наиболее компактные модели: они максимально доступны для обучения и развёртывания, демонстрируют ненулевое, но далёкое от потолка качество рассуждений и потому обладают наибольшим пространством для улучшения. В качестве базовой модели для всех экспериментов с дообучением выбрана \textbf{Qwen3-0.6B}~--- наименьшая модель семейства, помещающаяся на одну GPU класса H100 даже при сравнительно длинных rollout'ах и обучении с LoRA-адаптерами. Свойственная компактным моделям неустойчивость соблюдения формата ответа компенсируется включением соответствующих штрафов непосредственно в функцию награды (раздел~\ref{subsec:reward-design}), что позволяет использовать структурный сигнал даже при изначально низком качестве реконструкции графа. Для подхода~2 извлечения графов (раздел~\ref{subsec:graph-extraction}) в качестве вспомогательных моделей-экстракторов использовались Qwen3-8B и Llama-3.1-8B-Instruct, а в части экспериментов с Wikidata~--- Qwen3-32B.

\subsection{Дизайн функции графовой награды}
\label{subsec:reward-design}

Ключевой инновацией предлагаемого решения является использование структурного сходства графов рассуждений в качестве плотного сигнала вознаграждения. В отличие от классических подходов, где награда определяется только корректностью финального ответа, разработанная функция награды композитна и включает шесть аддитивных компонент, отражающих как формальную корректность ответа, так и качество цепочки рассуждений.

Пусть $\hat{y}$~--- сгенерированный моделью ответ, $\hat{G} = (\hat{V}, \hat{E})$~--- извлечённый из него граф рассуждений, $G^* = (V^*, E^*)$~--- эталонный граф из датасета, $y^*$~--- эталонная метка. Полная награда определяется как
\begin{equation}
\label{eq:reward-total}
R(\hat{y}, y^*, \hat{G}, G^*) = r_{\text{think}} + r_{\text{answer}} + r_{\text{parse}} + r_{\text{correct}} + \alpha \cdot r_{\text{graph}} + r_{\text{format}},
\end{equation}
где $\alpha \in [0, 1]$~--- масштабирующий коэффициент графовой компоненты (\texttt{graph\_coverage\_scale}), позволяющий плавно интерполировать между чисто результирующим обучением ($\alpha = 0$, baseline) и обучением с полной плотной графовой наградой ($\alpha = 1$).

Компоненты награды определяются следующим образом:
\begin{itemize}
    \setlength{\itemsep}{\smallskipamount}
    \item $r_{\text{think}}$~--- бинарный сигнал наличия корректно открытого и закрытого блока \texttt{<think>...</think>}: $+5$ при наличии, $-50$ при отсутствии хотя бы одного из тегов. Стимулирует модель явно проявлять цепочку рассуждений.
    \item $r_{\text{answer}}$~--- аналогичный сигнал для блока \texttt{<answer>...</answer>}: $+5$ или $-50$.
    \item $r_{\text{parse}} = +5$~--- бонус за успешный синтаксический разбор JSON-блока списка рёбер из \texttt{<answer>}; начисляется только если ответ удалось привести к корректному списку троек.
    \item $r_{\text{correct}} = +100$~--- крупное вознаграждение за корректность финальной метки задачи (\texttt{entailment}/\texttt{not entailment}). Сохраняет совместимость с классическим outcome-based RL.
    \item $r_{\text{graph}} = 100 \cdot \mathrm{Cov}(\hat{G}, G^*)$~--- основная графовая компонента, где $\mathrm{Cov}$~--- покрытие эталонного подграфа, вычисляемое процедурой \texttt{validate\_subgraph} согласно формуле~(\ref{eq:coverage}) из раздела~\ref{subsec:metrics-formalization}. Эта компонента предоставляет градации в зависимости от того, насколько корректно модель воспроизвела структуру вывода, а не только финальный ответ.
    \item $r_{\text{format}} = -\min(10 \cdot |c|, 100)$~--- штраф за наличие постороннего текста $c$ вне обязательных тегов \texttt{<think>} и \texttt{<answer>}, пропорциональный длине этого текста с ограничением сверху по модулю.
\end{itemize}

Структура композитной награды и вклад отдельных компонент схематически показаны на рис.~\ref{fig:reward}.

\begin{figure}[H]
\centering
\begin{tikzpicture}[
    comp/.style={draw, rounded corners, align=center, minimum height=8mm, minimum width=26mm, font=\footnotesize, fill=black!3},
    sumb/.style={draw, circle, minimum size=11mm, font=\large, fill=black!8},
    arr/.style={-{Stealth[length=2mm]}, thick},
    node distance=3mm,
]
\node[comp] (think) {$r_{\text{think}}$\\\scriptsize формат think};
\node[comp, below=of think] (answer) {$r_{\text{answer}}$\\\scriptsize формат answer};
\node[comp, below=of answer] (parse) {$r_{\text{parse}}$\\\scriptsize разбор JSON};
\node[comp, below=of parse] (correct) {$r_{\text{correct}}$\\\scriptsize верная метка};
\node[comp, below=of correct] (graph) {$\alpha\cdot r_{\text{graph}}$\\\scriptsize покрытие $G^*$};
\node[comp, below=of graph] (format) {$r_{\text{format}}$\\\scriptsize штраф за мусор};

\node[sumb, right=28mm of correct, yshift=4mm] (sum) {$\Sigma$};
\node[comp, right=14mm of sum, minimum width=22mm] (total) {$R$\\\scriptsize итоговая награда};

\foreach \n in {think,answer,parse,correct,graph,format}{
    \draw[arr] (\n.east) -- (sum.west);
}
\draw[arr] (sum) -- (total);
\node[font=\scriptsize, above=1mm of graph.north, align=center] {\textbf{графовая компонента}};
\end{tikzpicture}
\caption{Композитная функция награды: шесть аддитивных компонент. Коэффициент $\alpha$ управляет вкладом графовой компоненты ($\alpha=0$~--- baseline)}
\label{fig:reward}
\end{figure}

Такая декомпозиция награды решает сразу несколько задач: компоненты $r_{\text{think}}, r_{\text{answer}}, r_{\text{format}}$ обеспечивают воспроизводимый формат вывода, необходимый для дальнейшего автоматического анализа; $r_{\text{parse}}$ дополнительно стимулирует синтаксическую корректность графовой части; $r_{\text{correct}}$ сохраняет основной критерий задачи; $r_{\text{graph}}$~--- ключевой инновационный компонент~--- превращает разреженный бинарный сигнал в плотный, частично информирующий о качестве рассуждений даже при неверном финальном ответе. Введение коэффициента $\alpha$ позволяет в рамках единого фреймворка проводить контрастное сравнение: $\alpha = 0$ соответствует baseline-варианту, а $\alpha > 0$~--- предлагаемому методу.

\textbf{Награда для открытого графа знаний (Wikidata).} В экспериментах на датасетах, построенных по Wikidata, эталонный полный граф знаний недоступен в явном виде, а извлечённые моделью сущности и отношения требуют сопоставления с реальными узлами Wikidata. Для этого случая в \texttt{TunePPO} реализована альтернативная функция награды \texttt{GraphMultihopQAReward} (\texttt{ppotune/reward.py}), которая использует вспомогательный экстрактор (в проведённых экспериментах~--- Qwen3-32B) для приведения предсказаний к каноническому виду и сопоставление по нормализованному расстоянию Левенштейна с порогом $0{,}8$. Режим сопоставления (\texttt{graph\_mode = entity\_only}) учитывает только корректность извлечённых сущностей-ответов, что соответствует специфике вопросно-ответных задач над KG.

\subsection{Алгоритм дообучения}
\label{subsec:training-algorithm}

В качестве алгоритма дообучения с подкреплением выбран \textbf{GRPO} (Group Relative Policy Optimization)~\cite{shao2024deepseekmath}~--- модификация PPO, исключающая обучение отдельной value-сети за счёт нормализации преимуществ внутри группы из $G$ ответов на один и тот же запрос. Выбор GRPO обусловлен двумя свойствами, критичными для задачи: (1)~снижение требований к памяти примерно вдвое по сравнению с классическим PPO, что позволяет дообучать модели объёмом 0.6B параметров на одной GPU класса H100; (2)~естественная пригодность для разреженных и композитных наград, к которым относится разработанная функция~(\ref{eq:reward-total}).

\textbf{Конфигурация обучения.} По результатам предварительных экспериментов и анализа литературы по GRPO зафиксирована следующая конфигурация (значения приведены для основной серии экспериментов с Qwen3-0.6B на RuleTaker):
\begin{itemize}
    \setlength{\itemsep}{\smallskipamount}
    \item Размер группы при rollout $G = 8$ ответов на каждый запрос, эффективный размер батча запросов $64$; \texttt{ppo\_epochs} $= 2$ внутри одного шага обновления, \texttt{ppo\_batch\_size} $= 32$, \texttt{gradient\_accumulation\_steps} $= 4$.
    \item Коэффициент KL-регуляризации к опорной политике $\beta = 0{,}1$, параметр клиппинга $\varepsilon = 0{,}2$.
    \item Максимальная длина запроса $1024$ токена; максимальная длина ответа~--- до $4096$ токенов в основной серии на RuleTaker (платформа VERL) и $512$--$1024$ токенов для задач Wikidata (платформа torchtune) в зависимости от сложности рассуждения.
    \item Температура генерации $T = 0{,}7$, число шагов обучения $\approx 232$.
    \item LoRA-адаптеры~\cite{hu2022lora} ранга $r = 64$ и $\alpha = 16$, применяемые к матрицам внимания (\texttt{q\_proj}, \texttt{k\_proj}, \texttt{v\_proj}, \texttt{output\_proj}) и к слоям MLP. Полная заморозка базовых весов сокращает число обучаемых параметров примерно до $1\%$ от общего числа.
    \item Оптимизатор \texttt{PagedAdamW} из библиотеки \texttt{bitsandbytes} с learning rate $1 \cdot 10^{-4}$, обеспечивающий пейджинг состояния оптимизатора в системную память и тем самым дополнительно снижающий пиковое потребление VRAM.
\end{itemize}

\textbf{Платформы обучения.} В рамках работы дообучение реализовано параллельно на двух промышленных платформах: для основной серии экспериментов с датасетом RuleTaker использован фреймворк \textbf{VERL}~\cite{sheng2024verl}, обеспечивающий эффективную интеграцию vLLM-инференса в RL-цикл; для экспериментов на синтетических Wikidata-задачах применён собственный стек на базе \textbf{torchtune}, оформленный в виде пакета \texttt{TunePPO}. Такое дублирование позволяет независимо проверить устойчивость предлагаемого метода относительно деталей реализации алгоритма и одновременно подготовить инфраструктуру для дальнейших исследований.

\subsection{Разбор примера}
\label{subsec:worked-example}

Проиллюстрируем работу всего конвейера на конкретном двухшаговом примере RuleTaker.

\textbf{Вход.} Контекст содержит факт <<\texttt{cow chases lion}>> и два правила: <<если нечто преследует льва, то льву нужен тигр>> и <<если льву нужен тигр, то лев замёрз>>. Проверяемое утверждение (вопрос)~--- <<\texttt{lion is cold}>>. Корректный ответ~--- \texttt{entailment}, а эталонное отображение посылок имеет вид
\[
\Pi^* = \big\{\, \texttt{lion needs tiger} \mapsto [[\texttt{cow chases lion}]],\quad \texttt{lion is cold} \mapsto [[\texttt{lion needs tiger}]] \,\big\}.
\]

\textbf{Эталонный граф доказательства} показан на рис.~\ref{fig:example-graph}: факт-лист порождает промежуточное следствие, из которого выводится вопрос-корень.

\begin{figure}[H]
\centering
\begin{tikzpicture}[
    nd/.style={draw, rounded corners, align=center, font=\footnotesize, minimum height=8mm, inner sep=4pt},
    fact/.style={nd, fill=black!3},
    mid/.style={nd, fill=black!6},
    root/.style={nd, fill=black!10, thick},
    arr/.style={-{Stealth[length=2mm]}, thick},
]
\node[fact] (f) {\texttt{cow chases lion}\\\scriptsize (факт-лист)};
\node[mid, right=22mm of f] (m) {\texttt{lion needs tiger}\\\scriptsize (следствие)};
\node[root, right=22mm of m] (r) {\texttt{lion is cold}\\\scriptsize (вопрос-корень)};
\draw[arr] (f) -- (m);
\draw[arr] (m) -- (r);
\end{tikzpicture}
\caption{Эталонный граф доказательства для примера (двухшаговый вывод)}
\label{fig:example-graph}
\end{figure}

\textbf{Генерация модели.} Обученная политика выдаёт ответ в требуемом формате:
\begin{center}
\texttt{<think>cow chases lion is a fact; by rule 1 lion needs tiger; by rule 2 lion is cold</think>}\\
\texttt{<answer>[["cow chases lion","lion needs tiger"],["lion needs tiger","lion is cold"]]}\\
\texttt{entailment</answer>}
\end{center}

\textbf{Извлечение и валидация.} Парсер извлекает два ребра, формируя $\hat{G}$, изоморфный эталону. Процедура \texttt{validate\_subgraph} стартует от вопроса \texttt{lion is cold}: его посылка $\{\texttt{lion needs tiger}\}$ совпадает с эталонным вариантом~--- узел валиден; рекурсивно узел \texttt{lion needs tiger} с посылкой $\{\texttt{cow chases lion}\}$ также валиден; \texttt{cow chases lion}~--- факт-лист, валиден. Итого $\mathrm{Cov} = 3/3 = 1{,}0$.

\textbf{Расчёт награды.} Компоненты награды для данного ответа сведены в таблицу~\ref{tab:reward-example}.

\begin{table}[H]
\centering
\caption{Покомпонентный расчёт награды для разобранного примера}
\label{tab:reward-example}
\begin{tabular}{|l|c|l|}
\hline
\textbf{Компонента} & \textbf{Значение} & \textbf{Условие} \\
\hline
$r_{\text{think}}$   & $+5$   & корректный блок \texttt{<think>} \\
$r_{\text{answer}}$  & $+5$   & корректный блок \texttt{<answer>} \\
$r_{\text{parse}}$   & $+5$   & JSON-массив рёбер разобран \\
$r_{\text{correct}}$ & $+100$ & метка \texttt{entailment} верна \\
$\alpha\cdot r_{\text{graph}}$ & $+100$ & $\mathrm{Cov}=1{,}0$, $\alpha=1$ \\
$r_{\text{format}}$  & $0$    & постороннего текста нет \\
\hline
\textbf{Итого $R$}   & $\mathbf{+215}$ & \\
\hline
\end{tabular}
\end{table}

Для сравнения, в baseline-режиме ($\alpha = 0$) тот же ответ получил бы награду $115$: графовая компонента, поощряющая корректную \emph{структуру} вывода, в нём отсутствует.

\textbf{Отрицательный пример и частичное покрытие.} Рассмотрим два дополнительных случая, демонстрирующих градации графовой награды.
\begin{itemize}
    \setlength{\itemsep}{\smallskipamount}
    \item \textit{Корректный отказ от вывода.} Если эталонная метка~--- \texttt{not entailment} и модель выдала \emph{пустой} граф (ни одно заключение не выведено), это соответствует корректному отказу: полагается $\mathrm{Cov} = 1{,}0$, и графовая компонента максимальна. Если же при метке \texttt{not entailment} модель всё же построила рёбра, покрытие штрафуется до $\mathrm{Cov} = 0{,}5$.
    \item \textit{Частично корректная структура.} Пусть для заключения \texttt{lion needs tiger} модель указала посылку $\{\texttt{dog sees lion}\}$ вместо эталонной $\{\texttt{cow chases lion}\}$. Заключение присутствует в $G^*$, но набор посылок не совпадает ни с одним допустимым вариантом~--- узел признаётся \emph{частично валидным} со $\mathrm{score} = 0{,}5$. При наличии в графе одного полностью валидного и одного частично валидного узла покрытие составит $\mathrm{Cov} = (1 + 0{,}5)/2 = 0{,}75$, что транслируется в промежуточную графовую награду $\alpha \cdot 75$.
\end{itemize}
Именно такие градации отличают плотный графовый сигнал от разреженной бинарной награды: ответ с частично корректной цепочкой получает промежуточное поощрение, направляя обучение к структурно корректным рассуждениям.

\subsection{Полный конвейер построенного решения}
\label{subsec:overall-pipeline}

Совокупный конвейер дообучения, собирающий все описанные выше компоненты в единый процесс, состоит из следующих шагов:
\begin{enumerate}
    \setlength{\itemsep}{\smallskipamount}
    \item \textbf{Подготовка данных.} Исходный датасет (RuleTaker или сгенерированные Wikidata-задачи) проходит фильтрацию (исключение примеров с отрицаниями и из NatLang-конфигураций, схлопывание вершин-правил), разбиение на тренировочную и валидационную выборки в соотношении $0{,}9 : 0{,}1$ и сериализацию в JSON-формат VERL с полями \texttt{prompt}, \texttt{proof\_rules}, \texttt{label} (см.~раздел~\ref{subsec:dataset-construction}).
    \item \textbf{Инициализация политики.} Реализован этап холодного старта (cold-start SFT), адаптирующий базовую модель Qwen3-0.6B к формату структурированных рассуждений (подробно описан в разделе~\ref{subsec:impl-sft}). В основной серии \emph{контрастных} экспериментов, результаты которых приводятся в главе~\ref{sec:Chapter4}, в качестве начальной политики $\pi_0$ намеренно используется \emph{базовая} модель Qwen3-0.6B без SFT-инициализации: это изолирует вклад собственно графовой награды и обеспечивает корректность сравнения вариантов $\alpha = 0$ и $\alpha > 0$, обучаемых из одной и той же отправной точки. На матрицы внимания и MLP накладываются LoRA-адаптеры с указанной выше конфигурацией.
    \item \textbf{Rollout.} На каждом шаге обучения политика $\pi_\theta$ генерирует $G = 8$ ответов на каждый из запросов батча с температурой $T = 0{,}7$; генерация осуществляется через vLLM в формате \texttt{<think>...</think><answer>...</answer>}.
    \item \textbf{Извлечение графа.} Из каждого сгенерированного ответа парсится JSON-список рёбер из блока \texttt{<answer>}; при невозможности прямого разбора активируется fallback на \texttt{StructuredGraphCreator} (раздел~\ref{subsec:graph-extraction}).
    \item \textbf{Вычисление награды.} По каждому ответу вычисляются все шесть компонент награды~(\ref{eq:reward-total}); графовая компонента $r_{\text{graph}}$ опирается на покрытие эталонного подграфа~(\ref{eq:coverage}).
    \item \textbf{Обновление политики.} В рамках GRPO внутри каждой группы из $G$ ответов вычисляются нормализованные преимущества, после чего политика обновляется $\text{ppo\_epochs} = 2$ раза с клиппингом $\varepsilon = 0{,}2$ и KL-регуляризацией $\beta = 0{,}1$ к опорной политике $\pi_0$.
    \item \textbf{Валидация.} Каждые $K$ шагов на отложенной выборке рассчитываются метрики качества: точность по метке, $F_1$ по рёбрам, GES, а также детализированная статистика по каждой компоненте награды; результаты логируются в TensorBoard модулем \texttt{tensorboard\_logger.py}.
\end{enumerate}

Общая схема построенного конвейера приведена на рис.~\ref{fig:pipeline}.

\begin{figure}[H]
\centering
\begin{tikzpicture}[
    node distance=6mm and 10mm,
    box/.style={draw, rounded corners, align=center, minimum height=8mm, minimum width=30mm, font=\small, fill=black!3},
    io/.style={draw, align=center, minimum height=8mm, minimum width=30mm, font=\small, fill=black!8},
    arr/.style={-{Stealth[length=2mm]}, thick},
]
\node[io] (data) {Подготовка данных\\\footnotesize $(q, G^*)$, фильтрация, split};
\node[box, below=of data] (init) {Инициализация политики\\\footnotesize Qwen3-0.6B + LoRA};
\node[box, below=of init] (roll) {Rollout (vLLM)\\\footnotesize $G{=}8$ ответов, $T{=}0{,}7$};
\node[box, below=of roll] (extract) {Извлечение графа $\hat{G}$\\\footnotesize JSON-парсинг / fallback};
\node[box, below=of extract] (reward) {Вычисление награды\\\footnotesize 6 компонент, формула~(\ref{eq:reward-total})};
\node[box, below=of reward] (update) {Обновление политики (GRPO)\\\footnotesize нормировка преимуществ};
\node[io, below=of update] (val) {Валидация и логирование\\\footnotesize accuracy, Edge $F_1$, GES};

\draw[arr] (data) -- (init);
\draw[arr] (init) -- (roll);
\draw[arr] (roll) -- (extract);
\draw[arr] (extract) -- (reward);
\draw[arr] (reward) -- (update);
\draw[arr] (update) -- (val);
\draw[arr] (update.east) to[out=20, in=-20, looseness=1.8] node[right, font=\scriptsize, align=left] {следующий\\шаг} (roll.east);
\end{tikzpicture}
\caption{Общая схема конвейера дообучения с графовой наградой}
\label{fig:pipeline}
\end{figure}

Таким образом, построенное решение интегрирует разработанные формальные метрики качества цепочек рассуждений, процедуру извлечения графа из текстового ответа модели и композитную графовую функцию награды в единый GRPO-цикл, реализованный поверх двух промышленных платформ. Контрастное переключение между baseline ($\alpha = 0$) и предложенным методом ($\alpha > 0$) внутри одного фреймворка обеспечивает методологически корректное эмпирическое сравнение, результаты которого приведены в главе~\ref{sec:Chapter4}.